
\documentclass[11pt]{article}
\usepackage[margin=1in]{geometry}
\usepackage{amsmath,amssymb}
\usepackage{enumitem}
\usepackage[T1]{fontenc}
\usepackage{lmodern}
\setlength{\parskip}{0.5em}
\setlength{\parindent}{0pt}

\begin{document}

\begin{center}
{\LARGE \textbf{IB Physics 2016--2020}}\\
{\large \textbf{Mapped Collection for D4}}\\
\vspace{0.3em}
\textit{Faithful paraphrases of old-syllabus questions matched to induction, generators, power transmission, and closely related capacitor/energy-transfer material.}
\end{center}

\textbf{Mapping used.} Since 2016--2020 papers use the old syllabus, this file maps modern D4 mainly to old Topic 11 material (electromagnetic induction, power generation/transmission, capacitance) plus very closely related power-system questions.

\section*{Problems}

\begin{enumerate}[leftmargin=*, itemsep=1.2em]

\item \textbf{M17 TZ2 HL Paper 2, Q2(b) --- wave-driven generator}\\
A buoy floating in a vertical tube rises and falls with the sea surface. When it rises, a cable turns a generator on the sea bed; when it falls, the cable is rewound and no power is produced. The buoy's motion may be treated as SHM. The water wave has amplitude $4.3\,\text{m}$, wavelength $35\,\text{m}$, and speed $3.4\,\text{m s}^{-1}$.
\begin{enumerate}[label=(\alph*)]
\item Calculate the maximum vertical speed of the buoy.
\item Sketch the variation of generator output power with time over one cycle.
\item Explain why power is produced during only part of the oscillation.
\end{enumerate}

\item \textbf{M17 TZ2 HL Paper 2, Q6(b)--(d) --- ac power transmission in cables}\\
A cable made of many copper wires transfers energy from an ac generator to an electrical load. Each wire has length $35\,\text{km}$ and resistance $64\,\Omega$. The generator output power is $110\,\text{MW}$ and the peak potential difference is $150\,\text{kV}$. The resistivity of copper is $1.7\times 10^{-8}\,\Omega\text{m}$.
\begin{enumerate}[label=(\alph*)]
\item Calculate the radius of each wire.
\item Determine the peak current in the cable.
\item Calculate the power dissipated in the cable per unit length.
\item Two identical cables are then connected in parallel. Find the rms current in each cable.
\item Explain the variation of the magnetic force between the two parallel cables during an ac cycle.
\end{enumerate}

\item \textbf{M18 TZ2 HL Paper 2, Q7 --- pumped-storage hydroelectric system}\\
A pumped-storage hydroelectric station stores water behind a dam. Water descends through a vertical drop of $110\,\text{m}$ and drives a turbine-generator. The water flow rate is $1.2\times10^5\,\text{m}^3$ per minute and the density of water is $1.0\times10^3\,\text{kg m}^{-3}$.
\begin{enumerate}[label=(\alph*)]
\item Estimate the specific gravitational energy available from the stored water, with unit.
\item Show that the mean rate of decrease of gravitational potential energy is about $2.5\,\text{GW}$.
\item Explain how the operators can still make a profit even though energy is needed to pump water back uphill.
\end{enumerate}

\item \textbf{M19 TZ1 HL Paper 2, Q8(a) --- construction of a parallel-plate capacitor}\\
A student builds a capacitor of capacitance $68\,\text{nF}$ from aluminium foil and plastic film. The foil and film are $450\,\text{mm}$ wide; the film thickness is $55\,\mu\text{m}$; the film permittivity is $2.5\times10^{-11}\,\text{C}^2\text{N}^{-1}\text{m}^{-2}$.
\begin{enumerate}[label=(\alph*)]
\item Calculate the total length of foil required.
\item The film conducts if the electric field exceeds $1.5\,\text{MN C}^{-1}$. Calculate the maximum charge that may be stored safely.
\end{enumerate}

\item \textbf{M19 TZ1 HL Paper 2, Q8(b)--(c) --- capacitor discharge and induction}\\
The capacitor from the previous part is connected to a charging/discharging circuit with a battery of emf $12\,\text{V}$ and a resistor $R=1.2\,\text{k}\Omega$.
\begin{enumerate}[label=(\alph*)]
\item Calculate the time for the capacitor's charge to fall to $50\%$ of its fully charged value.
\item The ammeter in the discharge branch is replaced by a coil. Explain why an emf is induced in the coil during discharge.
\item Suggest one change to the discharge circuit, other than altering the coil itself, that would increase the maximum induced emf in the coil.
\end{enumerate}

\item \textbf{N19 HL Paper 2, Q9 --- capacitor network and thermal losses}\\
A charged capacitor or capacitor combination is connected in a new circuit and allowed to redistribute charge.
\begin{enumerate}[label=(\alph*)]
\item Calculate the initial energy stored using $E=\tfrac12QV$ or an equivalent expression.
\item Determine final charges on the capacitors after reconnection.
\item Find the total final electrical energy stored.
\item Explain why some of the original energy is dissipated as thermal energy in wires and resistors.
\end{enumerate}

\item \textbf{N20 HL Paper 2, Q9 --- rotating coil, induced emf, and step-up transmission}\\
A coil rotates in a magnetic field and is connected to a transmission system.
\begin{enumerate}[label=(\alph*)]
\item Explain, using magnetic flux linkage and Faraday's law, why an emf is induced in the rotating coil.
\item Given peak emf and current data, determine the rms voltage and rms current supplied to the load.
\item If transmission losses are reduced by using a step-up stage, determine the required factor by which the voltage must be increased.
\item The peak emf is doubled and the rotation period is halved. Sketch or describe the effect on the emf--time graph.
\end{enumerate}

\end{enumerate}

\end{document}
